\section{Case Study and Experimental Setup}
\label{sec:setup}

\subsection{The domain as a testbed, not as the scope}
\label{sec:setup-domain}

ChargeShield-FL is not an EV-charging system: it is an audit infrastructure
for federated training in OT environments, and charging is the scenario we
instrument it on. The choice is not arbitrary, because this domain satisfies
four conditions rarely found together.

A \emph{public dataset of real sessions} exists at useful scale, drawn from
genuinely distinct operational sites and carrying user identifiers, which makes
record-level membership a meaningful notion rather than a convention. The
operators are \emph{genuine commercial silos} rather than a synthetic partition
of one dataset, so cross-silo is a property of the data and not of the
experimental design. The infrastructure \emph{sits in a levelled environment}
under the Purdue model, which is what makes the observability gap of
\S\ref{sec:introduction} expressible at all. And session records are
\emph{simultaneously commercially sensitive and personal}, so the privacy
question is not hypothetical.

What is domain-specific in the implementation is a dataset adapter and a choice
of features. The ML Plane, the Privacy Auditor and the attack suite operate on
tensors and metadata and contain no notion of charging, vehicles or energy. We
exercise that independence with a second adapter for a structurally different
charging dataset, and report it as a property of the design rather than as an
experimental result, since one replication does not establish generality.

\subsection{Data}
\label{sec:setup-data}

We use ACN-Data~\cite{lee2019acndata} across its three real sites: Caltech
(54~EVSE, 31{,}404 sessions), JPL (50~EVSE, 33{,}629) and Office~1 (8~EVSE,
1{,}680), for 66{,}713 sessions in total.

\subsection{Why an autoencoder}
\label{sec:setup-model}

The detector is an autoencoder flagging a session whose reconstruction error
exceeds a threshold. Three reasons motivate this, and the third bears on
whether privacy can be measured at all.

\emph{Unsupervised out of necessity.} In a real OT setting attacks are rare,
heterogeneous and unlabelled: no operator holds enough annotated incidents to
train a classifier, which is the very reason federated learning is proposed
here. A model that learns normality and flags departures from it is the only
family applicable without shared labels.

\emph{Compatible with the deployment constraint.} The nodes running local
training are station controllers, not accelerators. A few-hundred-parameter
model is what that placement admits.

\emph{Directly auditable.} Reconstruction error is a per-record scalar, and it
is exactly the quantity loss-based membership attacks threshold on
(\S\ref{sec:attacks-strength}). An autoencoder is therefore auditable without
adaptation: no notion of confidence and no softmax output are required, and the
same quantity serves as the detector's anomaly score and the adversary's attack
score. That coincidence removes one degree of freedom between what the system
computes and what the attacker observes.

The model is deliberately conventional. We measure privacy risk on a detector
representative of what an operator would actually deploy, not on an
architecture chosen to make the measurement interesting; a novel architecture
would say less to anyone deciding on a deployment. Nothing in the framework
depends on this choice --- swapping the model requires an adapter, not a change
to the audit infrastructure.

\subsection{Training regime}
\label{sec:setup-training}

Autoencoder $6 \rightarrow 16 \rightarrow 8 \rightarrow 4 \rightarrow 8
\rightarrow 16 \rightarrow 6$, 570 parameters, over six continuous session
features: delivered energy, mean power, requested energy, declared
availability window, connection hour and duration. Adam, learning rate
$10^{-3}$, batch size 32, FedAvg with a proximal term.

\emph{Splits and pools.} Each client partitions its sessions into a local
training set and a held-out set. \emph{Members} are sessions in a client's
local training set; \emph{non-members} are held-out sessions from the same
site, never seen by any client. The evaluation pool is balanced by construction
and sampled once before the round loop, so every round scores the same samples.

\emph{Deliberately induced overfitting.} The local-epoch count is not a neutral
choice. It was set to induce local overfitting --- that is, to \emph{create}
the membership signal the experiment was then meant to measure --- and was
documented as such in this project's configuration from the outset. We state it
because it strengthens rather than weakens a null result: we actively tried to
produce the signal.

\emph{Repetitions.} Five independent seeds (42, 123, 456, 789, 1234). Each
re-initialises weights, the train/holdout split, shuffle order and every
downstream stochastic draw, so the runs are genuinely independent repetitions
rather than repeated readings of one run.

\subsection{What we do not measure: detection accuracy}
\label{sec:setup-notmeasured}

We do not report the model's detection accuracy, and the reason lies in the
data: ACN-Data carries no intrusion labels, and the field provided for them is
null across all 66{,}713 sessions. To our knowledge no public EV-charging
dataset supplies them.

This is the same scarcity that motivates federated learning in this domain. If
each operator held enough annotated incidents it would train a detector alone
and collaboration would be unnecessary; it is precisely because labels are
absent that an unsupervised model is trained jointly, and it is that
collaboration which raises the privacy question we study.

We therefore evaluate the detector for the one property our question requires
--- what its training reveals about the data that produced it --- and not for
its ability to detect attacks. We considered and rejected injecting synthetic
anomalies: that would measure the detector against the generator producing
them, and would not transfer. The model is representative in architecture,
scale and training regime of what an operator would deploy, while its
operational effectiveness remains outside this work's scope.

\subsection{The privacy unit does not match the subject}
\label{sec:setup-privacyunit}

Differential privacy bounds the influence of a data unit fixed by the chosen
neighbouring relation~\cite{dwork2006calibrating}. In our two mechanisms that
unit is the individual record, under per-example clipping, or a client's entire
contribution, under delta clipping. Neither coincides with the unit a person
would recognise as their own, and on ACN-Data the gap is measurable rather than
hypothetical.

A user identifier is present in 48{,}404 of 66{,}713 sessions, with markedly
uneven coverage across sites: $94\%$ at Caltech, $52\%$ at JPL, $35\%$ at
Office~1. Among identified sessions, 1{,}028 distinct users contribute a median
of \textbf{14 sessions each}, and 87 of them --- \textbf{$8.5\%$} --- charge at
more than one site, accounting for 7{,}618 sessions.

Two distinct gaps follow. Record granularity does not protect a person, because
14 records compose: budget spent on each accumulates, and the
individual-level guarantee is correspondingly weaker than the nominal
per-record one. Client granularity does not protect them either, because for
$8.5\%$ of users the contribution crosses organisational boundaries, and no
operator can clip data it does not hold.

\emph{A case observable within our own positive control.} Of the forty
templates drawn for the balanced control, five carry a user identifier, and
three belong to the same person: one session assigned to the member group and
two to the non-member group, all recorded at the same connection point over
five months. A record-level attack answers correctly the question it poses ---
was this session in the training set? --- yet for that driver a negative answer
on two records coexists with an affirmative one on a third. None of the three
answers tells the person what they would want to know, namely whether their
charging activity entered the model; the answer to that is yes. We report this
as an illustration rather than a statistic: with five identified templates out
of forty it is anecdotal by construction, and the measurement remains the
dataset-wide one above.

\emph{Why we do not evaluate subject-level DP.} The natural remedy bounds each
subject's aggregate contribution, which requires grouping gradients by person
rather than by record and raises sensitivity to the maximum number of sessions
per user. Beyond the implementation cost, evaluation on this data would be
partial: at $35\%$ coverage in Office~1, where the positive control is run, the
subject unit is undefined for most sessions. We report it as future work
motivated by a measurement, not as a neglected option.

\subsection{Realisation}
\label{sec:setup-realisation}

A single-process simulation and a five-container NVFLARE/Containerlab
deployment (server, caltech, jpl, office1, fl-admin), with the same unmodified
attack and analysis code run against both. The main results come from the
simulation; the containerised deployment is supplementary validation.
